% Generated from locale/tr/content/first-order-logic/syntax-and-semantics/formation-sequences.tex; do not hand-edit.


\olfileid[tr]{fol}{syn}{fseq}

\olsection{Oluşum Dizileri}

!!{formula}s için tümevarımsal bir tanım vermek ve tamamlayıcı bir
yöntem olarak !!{formula}s hakkında özellikleri tümevarımla ispatlamak zarif
ve etkili bir yaklaşımdır. Bununla birlikte, !!{formula}s söz konusu
olduğunda kuruluş sürecini daha aşağıdan yukarıya, adım adım ele almak da yararlı olabilir. Bunu
burada \emph{oluşum dizisi} kavramıyla yapacağız.
%
Terimler ile !!{formula}s bu yolla tanıtılırken tümevarımsal tanımlara
gerek olmadığını göstermek için önce bir~$\Lang L$ dilinden
alınmış gelişigüzel bir simge dizisi kavramını sunuyoruz.

\begin{defn}[Simge dizileri]
\ollabel{defn:string}
$\Lang L$ birinci dereceden bir dil olsun. Bir \emph{$\Lang L$-dizisi},
$\Lang L$'nin simgelerinden oluşan sonlu bir dizidir. $\Lang L$ dilinin
bağlamdan açıkça belli olduğu yerlerde bir $\Lang L$-dizisinden çoğu
kez yalnızca \emph{simge dizisi} diye söz ederiz.
\end{defn}

\begin{ex}
Her birinci dereceden $\Lang L$ dili için bütün
$\Lang L$-!!{formula}s birer $\Lang L$-dizisidir, fakat bunun tersi
geçerli değildir. Örneğin 
\[)(\Obj v_0\lif\lexists\]
bir $\Lang L$-dizisidir, ama bir $\Lang L$-!!{formula} değildir.
\end{ex}

\begin{defn}[Terimler için oluşum dizileri]
\ollabel{defn:fseq-trm}
Sonlu bir $\Lang L$-dizileri dizisi $\tuple{t_0,\dotsc,t_n}$, $t \ident
t_n$ ise ve her $i \leq n$ için ya $t_i$ bir !!a{variable} ya da bir
!!a{constant} ise veya $\Lang L$, $k$-yerli bir !!{function}~$f$
içeriyor ve $t_i \ident f(t_{m_1},\dotsc,t_{m_k})$ olacak biçimde
$m_1,\dotsc,m_k < i$ indisleri bulunuyorsa, $t$ terimi için bir
\emph{oluşum dizisi}dir. Ayrım yapmak gerektiğinde terimlere ilişkin
oluşum dizilerine \emph{terim oluşum dizileri} diyeceğiz.
\end{defn}

\begin{ex}
Şu dizi
\[
    \tuple{\Obj c_0, \Obj v_0, \Atom{\Obj f^2_0}{\Obj c_0, \Obj v_0}, \Atom{\Obj f^1_0}{\Atom{\Obj f^2_0}{\Obj c_0, \Obj v_0}}}
\]
$\Atom{\Obj f^1_0}{\Atom{\Obj f^2_0}{\Obj c_0, \Obj v_0}}$ terimi için
bir oluşum dizisidir; şu dizi de öyledir:
\[
    \tuple{\Obj v_0, \Obj c_0, \Atom{\Obj f^2_0}{\Obj c_0, \Obj v_0}, \Atom{\Obj f^1_0}{\Atom{\Obj f^2_0}{\Obj c_0, \Obj v_0}}}.
\]
\end{ex}

\begin{defn}[Formüller için oluşum dizileri]
\ollabel{defn:fseq-frm}
Sonlu bir $\Lang L$-dizileri dizisi $\tuple{!A_0,\dotsc,!A_n}$,
$!A \ident !A_n$ ise ve her $i \leq n$ için ya $!A_i$ atomik bir
!!{formula} ise ya da aşağıdakilerden biri geçerli olacak biçimde
$j,k < i$ indisleri ve bir !!a{variable}~$x$ varsa, $!A$ için bir
\emph{oluşum dizisi}dir:
\begin{enumerate}
    \tagitem{prvNot}{$!A_i \ident \lnot !A_j$.}{}%
    \tagitem{prvAnd}{$!A_i \ident (!A_j \land !A_k)$.}{}%
    \tagitem{prvOr}{$!A_i \ident (!A_j \lor !A_k)$.}{}%
    \tagitem{prvIf}{$!A_i \ident (!A_j \lif !A_k)$.}{}%
    \tagitem{prvIff}{$!A_i \ident (!A_j \liff !A_k)$.}{}%
    \tagitem{prvAll}{$!A_i \ident \lforall[x][!A_j]$.}{}%
    \tagitem{prvEx}{$!A_i \ident \lexists[x][!A_j]$.}{}%
\end{enumerate}
Ayrım yapmak gerektiğinde formüllere ilişkin oluşum dizilerine
\emph{formül oluşum dizileri} diyeceğiz.
\end{defn}

\begin{ex}
\[
\tuple{
    \Atom{\Obj A^1_0}{\Obj v_0},
    \Atom{\Obj A^1_1}{\Obj c_1},
    (\Atom{\Obj A^1_1}{\Obj c_1} \land \Atom{\Obj A^1_0}{\Obj v_0}),
    \lexists[\Obj v_0][(\Atom{\Obj A^1_1}{\Obj c_1} \land \Atom{\Obj A^1_0}{\Obj v_0})]
}
\]
şu formülün bir oluşum dizisidir:
$\lexists[\Obj v_0][(\Atom{\Obj A^1_1}{\Obj c_1} \land \Atom{\Obj
A^1_0}{\Obj v_0})]$. Aşağıdaki de bu formülün bir oluşum dizisidir:
\begin{multline*}
\tuple{
    \Atom{\Obj A^1_0}{\Obj v_0},
    \Atom{\Obj A^1_1}{\Obj c_1},
    (\Atom{\Obj A^1_1}{\Obj c_1} \land \Atom{\Obj A^1_0}{\Obj v_0}),
    \Atom{\Obj A^1_1}{\Obj c_1},\\
    \lforall[\Obj v_1][\Atom{\Obj A^1_0}{\Obj v_0}],
    \lexists[\Obj v_0][(\Atom{\Obj A^1_1}{\Obj c_1} \land \Atom{\Obj A^1_0}{\Obj v_0})]
}.
\end{multline*}
%
İkinci örnekten görüldüğü gibi oluşum dizileri, gereksiz olan veya
kuruluşa katkıda bulunmayan !!{formula}s, yani ``artıklar'' içerebilir.
\end{ex}

\begin{prop}\ollabel{prop:formed}
$\Frm[L]$ içindeki her !!{formula}~$!A$ bir oluşum dizisine sahiptir.
\end{prop}

\begin{proof}
$!A$ atomik olsun. O zaman $\tuple{!A}$ dizisi, $!A$ için bir oluşum
dizisidir.
%
Şimdi $!B$ ile~$!C$'nin sırasıyla $\tuple{!B_0,\dotsc,!B_n}$ ve
$\tuple{!C_0,\dotsc,!C_m}$ oluşum dizilerine sahip olduğunu varsayalım.
%
\begin{enumerate}
  \tagitem{prvNot}{$!A \ident \lnot !B$ ise
      $\tuple{!B_0,\dotsc,!B_n,\lnot !B_n}$,
      $!A$ için bir oluşum dizisidir.}{}
  \tagitem{prvAnd}{$!A \ident (!B \land !C)$ ise
      $\tuple{!B_0,\dotsc,!B_n,!C_0,\dotsc,!C_m,(!B_n \land !C_m)}$,
      $!A$ için bir oluşum dizisidir.}{}
  \tagitem{prvOr}{$!A \ident (!B \lor !C)$ ise
      $\tuple{!B_0,\dotsc,!B_n,!C_0,\dotsc,!C_m,(!B_n \lor !C_m)}$,
      $!A$ için bir oluşum dizisidir.}{}
  \tagitem{prvIf}{$!A \ident (!B \lif !C)$ ise
      $\tuple{!B_0,\dotsc,!B_n,!C_0,\dotsc,!C_m,(!B_n \lif !C_m)}$,
      $!A$ için bir oluşum dizisidir.}{}
  \tagitem{prvIff}{$!A \ident (!B \liff !C)$ ise
      $\tuple{!B_0,\dotsc,!B_n,!C_0,\dotsc,!C_m,(!B_n \liff !C_m)}$,
      $!A$ için bir oluşum dizisidir.}{}
  \tagitem{prvAll}{$!A \ident \lforall[x][!B]$ ise
      $\tuple{!B_0,\dotsc,!B_n,\lforall[x][!B_n]}$,
      $!A$ için bir oluşum dizisidir.}{}
  \tagitem{prvEx}{$!A \ident \lexists[x][!B]$ ise
      $\tuple{!B_0,\dotsc,!B_n,\lexists[x][!B_n]}$,
      $!A$ için bir oluşum dizisidir.}{}
\end{enumerate}
!!{formula}s üzerinde tümevarım ilkesi uyarınca her !!{formula} bir
oluşum dizisine sahiptir.
\end{proof}

Tersini de ispatlayabiliriz. Bu önemlidir; çünkü formülleri tanımlamanın
iki yolunun eşdeğer olduğunu, yani aynı sonuçları verdiğini gösterir.
Ayrıca oluşum dizilerinin uzunluğu üzerinde olağan tümevarım kullanarak
formüller hakkında teoremler ispatlayabileceğimiz anlamına gelir.

\begin{lem}
\ollabel{lem:fseq-init}
$\tuple{!A_0,\dotsc,!A_n}$, $!A_n$ için bir oluşum dizisi ve $k \leq n$
olsun. O zaman $\tuple{!A_0,\dotsc,!A_k}$, $!A_k$ için bir oluşum
dizisidir.
\end{lem}

\begin{proof}
Alıştırma.
\end{proof}

\begin{prob}
\olref[fol][syn][fseq]{lem:fseq-init}'i ispatlayın.
\end{prob}

\begin{thm}
\ollabel{thm:fseq-frm-equiv}
$\Frm[L]$, $!A$ için bir formül oluşum dizisi bulunan bütün $\Lang
L$-dizileri~$!A$ dizilerinin kümesidir.
\end{thm}

\begin{proof}
$F$, $\Lang L$ dilindeki simgelerden oluşan ve bir oluşum dizisine sahip
bütün simge dizilerinin kümesi olsun. \olref[fol][syn][fseq]{prop:formed}'da
$\Frm[L] \subseteq F$ olduğunu gördük; şimdi ters kapsamayı ispatlayalım.

$!A$'nın bir $\tuple{!A_0,\dotsc,!A_n}$ oluşum dizisi olduğunu
varsayalım. $n$ üzerinde kuvvetli tümevarımla $!A \in \Frm[L]$ olduğunu
ispatlarız. Tümevarım varsayımımız, son indisi $m < n$ olan bir oluşum
dizisine sahip her simge dizisinin $\Frm[L]$'de bulunduğudur. Oluşum
dizisi tanımı uyarınca ya $!A \ident !A_n$ atomiktir ya da aşağıdaki
durumlardan birinin geçerli olmasını sağlayan $j,k < n$ indisleri vardır:
\begin{enumerate}
    \tagitem{prvNot}{$!A \ident \lnot !A_j$.}{}%
    \tagitem{prvAnd}{$!A \ident (!A_j \land !A_k)$.}{}%
    \tagitem{prvOr}{$!A \ident (!A_j \lor !A_k)$.}{}%
    \tagitem{prvIf}{$!A \ident (!A_j \lif !A_k)$.}{}%
    \tagitem{prvIff}{$!A \ident (!A_j \liff !A_k)$.}{}%
    \tagitem{prvAll}{$!A \ident \lforall[x][!A_j]$.}{}%
    \tagitem{prvEx}{$!A \ident \lexists[x][!A_j]$.}{}%
\end{enumerate}
Şimdi durumlara göre akıl yürütelim. $!A$ atomikse
$!A_n \in \Frm[L]$'dir. Bunun yerine
$!A \ident (!A_j \land !A_k)$ olduğunu varsayalım.
\olref[fol][syn][fseq]{lem:fseq-init} uyarınca
$\tuple{!A_0,\dotsc,!A_j}$ ve $\tuple{!A_0,\dotsc,!A_k}$ sırasıyla
$!A_j$ ile~$!A_k$ için oluşum dizileridir. Bunlar $!A$'nın oluşum
dizisinin öz başlangıç alt dizileridir ve son indisleri $n$'den küçüktür.
Dolayısıyla tümevarım varsayımı uyarınca $!A_j$ ile~$!A_k$,
$\Frm[L]$'dedir; bir !!a{formula} tanımı uyarınca
$(!A_j \land !A_k)$ de öyledir. Öteki durumlar koşut akıl yürütmelerle
elde edilir.
\end{proof}

Terimlere ilişkin oluşum dizileri, !!{formula}s için olanlara benzer
özelliklere sahiptir.

\begin{prop}
\ollabel{prop:fseq-trm-equiv}
$\Trm[L]$, $t$ için bir terim oluşum dizisi bulunan bütün $\Lang L$-dizileri
$t$ dizilerinin kümesidir.
\end{prop}

\begin{proof}
Alıştırma.
\end{proof}

\begin{prob}
\olref[fol][syn][fseq]{prop:fseq-trm-equiv}'i ispatlayın.
İpucu: \olref[fol][syn][fseq]{thm:fseq-frm-equiv}'in ispatında kullanılan
yönteme benzer bir yöntem kullanın.
\end{prob}

Oluşum dizilerinde iki tür ``artık'' görülebilir: yinelenen öğeler ve
formülün ya da terimin kuruluşuyla ilgisiz öğeler. En küçük oluşum
dizilerine bakarak ikisini de eleyebiliriz.

\begin{defn}[En küçük oluşum dizileri]
\ollabel{defn:minimal-fseq}
Bir $!A$ formülü için $\tuple{!A_0, \ldots, !A_n}$ oluşum dizisi,
$!A$ için başka her oluşum dizisi~$s$ için $s$'nin uzunluğu $n+1$'den büyük ya da
ona eşitse, $!A$ için bir \emph{en küçük oluşum dizisi}dir.

Benzer biçimde bir $t$ terimi için $\tuple{t_0, \ldots, t_n}$ oluşum
dizisi, $t$ için başka her oluşum dizisi~$s$ için $s$'nin uzunluğu $n+1$'den büyük
ya da ona eşitse, $t$ için bir \emph{en küçük oluşum dizisi}dir.
\end{defn}

Bir formülün veya terimin birden çok en küçük oluşum dizisi olabileceğine
dikkat edin; fakat bunlar tam olarak aynı simge dizilerini içerir.

\begin{prop}
\ollabel{prop:subformula-equivs}
Aşağıdakiler birbirine eşdeğerdir:
\begin{enumerate}
    \item $!B$, $!A$'nın bir alt !!{formula}{}dür.
    \item $!B$, $!A$'nın her oluşum dizisinde geçer.
    \item $!B$, $!A$'nın bir en küçük oluşum dizisinde geçer.
\end{enumerate}
\end{prop}

\begin{proof}
Alıştırma.
\end{proof}

\begin{prob}
\olref[fol][syn][fseq]{prop:subformula-equivs}'i ispatlayın.
\end{prob}

\begin{history}
Oluşum dizilerini Raymond Smullyan, \emph{First-Order Logic}
\citep{Smullyan1968} adlı ders kitabında ortaya koymuştur. Oluşum
dizilerinin başka özelliklerini \citet{Zuckerman1973} saptamıştır.
\end{history}

